\chapter{代码架构与关键实现}\label{app:code}

本附录介绍统一PINN期权定价系统的代码架构和关键模块的实现细节，供读者复现实验结果参考。

\section{项目目录结构}

系统代码组织如下：

\begin{verbatim}
uni_pinn/
├── README.md                    # 复现指南
├── option_pinn/
│   ├── unified_pinn_v2.py       # 统一PINN核心模型（v16）
│   ├── ref_solvers.py           # 参考解（BSM/CEV/Heston解析/半解析）
│   ├── eval_all.py              # 合成数据与市场数据评估脚本
│   ├── llm_test.py              # LLM路由层测试（20个测试用例）
│   ├── ablation.py              # 消融实验（三个消融变体）
│   ├── greeks_validation.py     # Greeks计算与验证
│   ├── generate_figures.py      # 论文图表生成
│   ├── llm_router.py            # LLM路由层（DeepSeek API）
│   ├── independent/             # 独立PINN基线
│   │   ├── bsm_pinn.py          # BSM门控网络（Dhiman & Hu 2023）
│   │   ├── cev_pinn.py          # CEV门控网络
│   │   ├── heston_hainaut.py    # Hainaut & Casas (2024) 参数化PINN
│   │   └── train_independent.py # 独立基线训练脚本
│   ├── data/
│   │   └── spy_quotedata.csv    # SPY期权链数据（2026-05-11）
│   └── results/                 # 评估输出（CSV，模型权重.pt）
└── thesis/                      # LaTeX论文源码
\end{verbatim}

\section{统一PINN核心模型}

\subsection{网络定义伪代码}

\begin{verbatim}
class UnifiedPINN(nn.Module):
    输入维度: 9 (S/Smax, v/vmax, t/T, sigma, beta,
                  kappa, theta, xi, rho)
    网络结构: 6层全连接, 每层128神经元, tanh激活
    最后一层: 权重初始化 N(0, 0.001^2), 偏置=0

    前向传播 forward(S, v, t, lambda):
        # 1. 归一化输入
        x = [S/300, v/1, t/T, sigma, beta,
             kappa, theta, xi, rho]

        # 2. 计算软掩码
        mask = tanh(xi / 0.05)^2

        # 3. 计算等效波动率
        sigma_eff = (1-mask)*sigma + mask*sqrt(v)

        # 4. 计算BS基线（解析公式）
        V_BS = black_scholes(S, K, tau, r, sigma_eff)

        # 5. 网络修正量
        net_out = MLP(x)  # 6层全连接

        # 6. 加法参数化输出
        V_hat = V_BS + K * net_out
        return V_hat
\end{verbatim}

\subsection{统一PDE残差计算伪代码}

\begin{verbatim}
def compute_pde_residual(model, S, v, t, lambda):
    # 使用自动微分计算偏导数
    V = model(S, v, t, lambda)

    dV_dt  = autograd(V, t)
    dV_dS  = autograd(V, S)
    d2V_dS2 = autograd(dV_dS, S)
    dV_dv  = autograd(V, v)
    d2V_dv2 = autograd(dV_dv, v)
    d2V_dSdv = autograd(dV_dS, v)

    # 软掩码
    mask = tanh(xi / 0.05)^2

    # 统一系数
    a = (1-mask)*sigma^2*S^(2*beta) + mask*v*S^2
    b = mask * rho * xi * v * S
    c = mask * xi^2 * v
    d = mask * kappa * (theta - v)

    # 统一PDE算子
    F = (dV_dt
         + 0.5*a*d2V_dS2
         + b*d2V_dSdv
         + 0.5*c*d2V_dv2
         + r*S*dV_dS
         + d*dV_dv
         - r*V)

    # 相对归一化残差
    residual = F / (|V| + 1)
    return mean(residual^2)
\end{verbatim}

\section{LLM路由层实现}

\subsection{系统提示词}

LLM路由层使用以下系统提示词（英文，以提高LLM的理解准确率）：

\begin{verbatim}
You are an option pricing parameter extractor.
Extract parameters from user input and output JSON.

Rules for model selection:
- If user mentions "Heston", "stochastic volatility",
  "vol-of-vol", kappa/theta/xi/rho: use "heston"
- If user mentions "CEV", "constant elasticity",
  "beta": use "cev"
- Otherwise: use "bsm"

Default values if not specified:
S=100, K=100, T=1.0, r=0.05, sigma=0.2,
beta=0.5, v0=0.04, kappa=2.0, theta=0.04,
xi=0.3, rho=-0.7

Output format (JSON only, no explanation):
{
  "model": "bsm"|"cev"|"heston",
  "option_type": "call"|"put",
  "S": float, "K": float, "T": float, "r": float,
  "sigma": float, "beta": float,
  "v0": float, "kappa": float, "theta": float,
  "xi": float, "rho": float
}
\end{verbatim}

\subsection{参数校验伪代码}

\begin{verbatim}
def validate_params(params):
    # 范围校验与截断
    params["sigma"] = clip(params["sigma"], 0.01, 2.0)
    params["beta"]  = clip(params["beta"],  0.1,  1.0)
    params["kappa"] = clip(params["kappa"], 0.1, 10.0)
    params["theta"] = clip(params["theta"], 0.001, 1.0)
    params["xi"]    = clip(params["xi"],    0.01,  2.0)
    params["rho"]   = clip(params["rho"],  -0.99,  0.99)
    params["v0"]    = clip(params["v0"],    0.001, 1.0)
    params["T"]     = clip(params["T"],     0.01, 10.0)
    params["S"]     = clip(params["S"],     1.0, 1000.0)
    params["K"]     = clip(params["K"],     1.0, 1000.0)
    return params
\end{verbatim}

\section{参考解实现}

\subsection{BSM解析解}

\begin{verbatim}
def black_scholes(S, K, tau, r, sigma, option_type):
    if tau <= 0:
        if option_type == "call":
            return max(S - K, 0)
        else:
            return max(K - S, 0)

    d1 = (log(S/K) + (r + 0.5*sigma^2)*tau) \
         / (sigma * sqrt(tau))
    d2 = d1 - sigma * sqrt(tau)

    if option_type == "call":
        return S*N(d1) - K*exp(-r*tau)*N(d2)
    else:
        return K*exp(-r*tau)*N(-d2) - S*N(-d1)
\end{verbatim}

\subsection{Heston半解析解（特征函数方法）}

Heston模型的欧式看涨期权价格通过特征函数方法计算：
\begin{equation}
  V_\text{Heston} = S\cdot P_1 - Ke^{-rT}\cdot P_2,
\end{equation}
其中$P_1$和$P_2$通过Gil-Pelaez反演公式计算：
\begin{equation}
  P_j = \frac{1}{2} + \frac{1}{\pi}\int_0^\infty
    \text{Re}\!\left[\frac{e^{-i\phi\ln K}\,\varphi_j(\phi)}{i\phi}\right]d\phi,
    \quad j=1,2,
\end{equation}
其中$\varphi_j(\phi)$为Heston模型的特征函数，具体形式见Heston（1993）原文。数值积分使用高斯-勒让德求积，积分上限截断为500，节点数为200，精度约$10^{-4}$。

\section{训练配置}

表~\ref{tab:train_config}列出了完整的训练配置，供复现实验参考。

\begin{table}[htbp]
  \centering
  \caption{统一PINN完整训练配置}
  \label{tab:train_config}
  \begin{tabular}{ll}
    \hline
    配置项 & 取值 \\
    \hline
    随机种子 & 42 \\
    训练设备 & NVIDIA GPU（CUDA）\\
    配点数$N_c$ & 每模型512点（BSM/CEV/Heston各512）\\
    边界点数$N_b$ & 每模型500点 \\
    数据锚点数$N_d$ & 每模型20点 \\
    优化器 & Adam（$\beta_1=0.9$，$\beta_2=0.999$）\\
    初始学习率 & $1\times 10^{-3}$ \\
    学习率调度 & 余弦退火（$\eta_\text{min}=1\times10^{-5}$）\\
    总训练步数 & 30000 \\
    损失权重$(w_\text{pde},w_\text{bc},w_\text{data})$ & $(1, 10, 100)$ \\
    软掩码阈值 & 0.05 \\
    最后一层初始化 & $\mathcal{N}(0, 0.001^2)$ \\
    BSM参数范围（训练） & $\sigma\in[0.10,0.35]$，共6个变体 \\
    CEV参数范围（训练） & $\sigma\in[0.10,0.35]$，$\beta\in[0.3,0.9]$，共12个变体 \\
    Heston参数范围（训练） & $\kappa\in[1.0,3.0]$，$\theta\in[0.02,0.08]$，\\
                           & $\xi\in[0.2,0.4]$，$\rho\in[-0.7,-0.5]$，共36个变体 \\
    \hline
  \end{tabular}
\end{table}
